\documentclass{article}

\usepackage{graphicx}
\usepackage{amsmath, amsthm, amssymb}
\usepackage{enumitem}
\usepackage{pgfplots}
\pgfplotsset{compat=1.18} % Ensures backward compatibility
\usepackage{mdframed}
\usepackage{float}
\usepackage{xcolor}
\usepackage{tcolorbox}
\usepackage{xurl}
\usepackage{colortbl}
\usepackage{hyperref} % Load this last in your preamble
\hypersetup{
    colorlinks=true,
    urlcolor=blue!70!black,
    linkcolor=purple!70!black
}
\usepackage[capitalise, noabbrev, nameinlink]{cleveref}

\theoremstyle{definition}
\newtheorem{definition}{Definition}
\newtheorem{example}{Example}
\surroundwithmdframed[
    hidealllines=true,
    leftline=true,
    linewidth=1.5pt,
    linecolor=gray!50,
    innerleftmargin=10pt,
    innerrightmargin=0pt,
    innertopmargin=2pt,
    innerbottommargin=2pt,
    skipabove=10pt,
    skipbelow=10pt
]{example}
\newtheorem{activity}{Activity}
\renewcommand{\theactivity}{\Alph{activity}}

\newcommand{\doenetlink}[2]{%
    \begin{tcolorbox}[
        colback=gray!6,
        colframe=blue!60!black,
        boxrule=0pt,
        leftrule=3pt,
        arc=2pt,
        left=8pt, right=8pt, top=6pt, bottom=6pt,
        before skip=10pt, after skip=10pt
    ]
    \textbf{\href{#2}{#1}} \\
    {\small\ttfamily \url{#2}}
    \end{tcolorbox}%
}
\newmdenv[
    backgroundcolor=gray!15,
    linecolor=black,
    linewidth=1pt,
    innertopmargin=10pt,
    innerbottommargin=10pt,
    userdefinedwidth=0.8\textwidth,
    align=center,
    nobreak=true,
    settings={\refstepcounter{calloutbox}},
    frametitle={\textbf{Box \thecalloutbox}}
]{calloutbox}
\newcounter{calloutbox}
\renewcommand{\thecalloutbox}{\arabic{calloutbox}}
\crefname{calloutbox}{box}{boxes}
\Crefname{calloutbox}{Box}{Boxes}

% Augmented matrix with a continuous light-gray vertical bar
% Usage: \begin{augmatrix}{ccc} ... \end{augmatrix} (specify coefficient columns only)
\newenvironment{augmatrix}[1]{%
  \left[\begin{array}{@{}#1!{\color{gray!75}\vrule}r@{}}%
}{%
  \end{array}\right]%
}

\newenvironment{rowops}{%
  \;\begin{array}{@{}l@{}}%
}{%
  \end{array}%
}


\title{Sequences}
\author{}
\date{}

\begin{document}

\maketitle

\section{Introduction to Sequences}\label{sec:sequences-intro}

% Define a sequence as a function whose domain is the positive integers (or non-negative integers).
% Introduce notation: a_n, \{a_n\}, terms of a sequence, index n.
% Calloutbox: Definition of a Sequence.
% Example: Writing out the first few terms from an explicit formula (including alternating signs (-1)^n and factorials n!).
% Example: Finding an explicit formula a_n from a list of initial terms.
% \doenetlink{Sequences Preview Activity}{https://doenet.org/code/...}

\section{Recursive Sequences}\label{sec:recursive-sequences}

% Introduce recursive definitions (initial condition(s) + recurrence relation).
% Calloutbox: Recursive Definition of a Sequence.
% Example: Computing terms of a recursively defined sequence (including Fibonacci-type two-term recurrence).

\section{Arithmetic and Geometric Sequences}\label{sec:arith-geom-sequences}

% Calloutbox: Arithmetic Sequences (common difference d, recursive a_n = a_{n-1} + d, explicit a_n = a_1 + (n-1)d).
% Example: Identifying arithmetic sequences, finding d, and writing explicit/recursive formulas.
% Calloutbox: Geometric Sequences (common ratio r, recursive a_n = r a_{n-1}, explicit a_n = a_1 r^{n-1}).
% Example: Identifying geometric sequences, finding r, and writing explicit/recursive formulas.

\section{Sigma Notation and Partial Sums}\label{sec:sigma-notation}

% Introduce summation (sigma) notation \sum_{k=1}^{n} a_k and nth partial sums S_n.
% Calloutbox: Summation Notation and Properties.
% Example: Evaluating sums and rewriting sums in sigma notation.

\end{document}
